\subsection{Linear Subspace Generated by a Subset}

\begin{tcolorbox}[title=Problem 1, breakable]
    If $T : V \longrightarrow W$ is linear and $A$
    is a subset of $V$ such that $T(A)$ is generating 
    for $W$, then $T$ is surjective.
\end{tcolorbox}

\begin{proof}
    Let $w$ be an arbitrary element in $W$.
    Since $T(A)$ generates $W$, $w$ is a linear combination of elements in $T(A)$.
    Suppose 
    \[
    w = c_1 a_1 + \cdots + c_n a_n \text{ such that } a_1, \ldots, a_n \in T(A).
    \]
    But $a_i = T(v_i)$ for some $v_i \in V$, so 
    \[
    c_1 a_1 + \cdots + c_n a_n = c_1 T(v_1) + \cdots + c_n T(v_n) = T(c_1 v_1) + \cdots + T(c_n v_n) = T(c_1 v_1 + \cdots + c_n v_n).
    \]
    The last two equalities hold since $T$ is linear.
    Thus $T$ is surjective.
\end{proof}

\begin{tcolorbox}[title=Problem 2, breakable]
    If $M$ and $N$ are linear subspaces of $V$ such that $V = M + N$ ($1.6.9$)
    and $T : V \longrightarrow V / N$ is the quotient mapping, then 
    $T(M) = V / N$ but $M$ need not equal $V$; thus we cannot conclude in 
    Exercise 1 that $A$ is a generating for $V$. Can you exhibit an example
    of such $V, M$ and $N$?
\end{tcolorbox}

\begin{proof}
    Consider $V = \mathbb{R}$, $M = \{0\}$, and $N = \mathbb{R}$. Then $V = M + N$, but $M \ne V$.
    Let $T : V \to V/N$ be the quotient map. Since $V/N = \{0\}$, we have $T(M) = \{0\} = V/N$.
\end{proof}

\begin{tcolorbox}[title=Problem 3, breakable]
    (Substitution Principle) Suppose $x, x_1, \ldots, x_n$ are 
    vectors such that $x$ is a linear combination of $x_1, \ldots, x_n$
    but not of $x_1, \ldots, x_{n - 1}$. Prove:
    \[[\{x_1, \ldots, x_n\}] = [\{x_1, \ldots, x_{n - 1}, x\}].\]
\end{tcolorbox}

\begin{proof}
    Since $x$ is a linear combination of $x_1, \ldots, x_n$ suppose 
    \[
    x = c_1 x_1 +\cdots + c_n x_n.
    \]
    Since $x$ is not a linear combination of $x_1, \ldots, x_{n-1}$, we have $c_n \neq 0$, so
    \[
    x_n = \frac{1}{c_n}\big(x - c_1 x_1 - \cdots - c_{n-1} x_{n-1}\big).
    \]
    Let $v \in [\{x_1, \ldots, x_n\}]$ and suppose 
    \[
    v = d_1 x_1 + \cdots + d_n x_n.
    \]
    Substituting for $x_n$, we get
    \[
    v = d_1 x_1 + \cdots + d_{n-1} x_{n-1} + d_n \frac{1}{c_n}\big(x - c_1 x_1 - \cdots - c_{n-1} x_{n-1}\big)
    \in [\{x_1, \ldots, x_{n-1}, x\}].
    \]
    Conversely, let $v \in  [\{x_1, \ldots, x_{n-1}, x\}]$ and suppose 
    \[
    v = d_1 x_1 + \cdots + d_{n - 1} x_{n - 1} + d x.
    \]
    Then 
    \[
    x = c_1 x_1 +\cdots + c_n x_n \implies dx = (d c_1) x_1 + \cdots + (d c_n) x_n.
    \]
    Thus 
    \[
    v = d_1 x_1 + \cdots + d_{n - 1} x_{n - 1} + (d c_1) x_1 + \cdots + (d c_n) x_n \in [\{x_1, \ldots, x_n\}].
    \]
\end{proof}

\begin{tcolorbox}[title=Problem 4, breakable]
    Let $S : V \longrightarrow W$ and $T : V \longrightarrow W$
    be linear mappings, and let $A$ be a subset of $V$ such that $V = [A]$.
    Prove that if $Sx = Tx$ for all $x \in A$, then $S = T$. (Informally,
    a lienar mapping is determined by its values on a generating set.)
\end{tcolorbox}

\begin{proof}
    Suppose $Sx = Tx$ for all $x \in A$.
    Let $x' \in V$ and since $V = [A]$ so 
    \[
    x' = c_1 a_1 + \cdots + c_n a_n.
    \]
    Then 
    \begin{align*}
        T(x') &= T(c_1 a_1 + \cdots + c_n a_n) \\
        &= T(c_1 a_1) + \cdots + T(c_n a_n) \\
        &= c_1 T(a_1) + \cdots + c_n T(a_n) \\
        &= c_1 S(a_1) + \cdots + c_n S(a_n) \\
        &= S(c_1 a_1) + \cdots + S(c_n a_n) \\
        &= S(c_1 a_1 + \cdots + c_n a_n) \\
        &= S(x').
    \end{align*}
    Thus $Sx' = Tx'$ for all $x' \in V$, so $S = T$.
\end{proof}

\begin{tcolorbox}[title=Problem 5, breakable]
    Let $V$ be a vector space, $A$ a nonempty subset of $V$. For each finite
    subset $\{x_1, \ldots, x_n\}$ of $A$, let 
    \[M(x_1, \ldots, x_n) = \{a_1 x_1 + \cdots + a_n x_n \mid a_1, \ldots, a_n \text{ scalars.}\};\]
    as noted in the proof of Corollary $3.1.5$, this is a linear subspace of $V$.
    Let $\mathcal{N}$ be the set of all linear subspaces of $V$ obtained in this way
    (by considering all possible finite subsets of $A$); by definition,
    $[A] = \bigcup \mathcal{N}$ (the union of the subspaces belonging to $\mathcal{N}$).
    Prove that $[A]$ is a linear subspace of $V$ by verifying criterion 1.6, Exercise 9.
\end{tcolorbox}

\begin{tcolorbox}[title=Problem 6, breakable]
    Let $V$ and $W$ be vector spaces, $V \times W$ the product vector space ($1.3.11$);
    let $A \subset V$ and $B \subset W$. Prove:
    \begin{enumerate}
        \item $[A \times B] \subset [A] \times [B]$,
        \item In general, $[A \times B] \ne [A] \times [B]$, {Hint: Consider $A = \emptyset, B = W$; or $V = W = R$ and $A = B = \{1\}$.}
        \item $[(A \times \{\emptyset\}) \cup (\{\emptyset\} \times B)] = [A] \times [B]$. {Hint 
        \[\left( \sum_{i = 1}^{m} c_i x_i, \sum_{i = 1}^{m} c_i (x_i, 0) + \sum_{j = 1}^{n} d_j (0, y_j) \right).\]    
        }
    \end{enumerate}
\end{tcolorbox}